\section{Discussion}

TE-KG addresses a practical integration problem in human TE research: evidence about the same named element is distributed across family references, genomic annotations, expression matrices and biomedical papers. The resource links these layers through shared TE identities and connected access workflows. Its main contribution is not the use of a graph in isolation. Rather, it is the combination of a PMID-bearing literature graph with classification, representative sequence and genomic records, three expression contexts and context-specific co-expression views, while keeping the interpretation boundary of each layer visible.

This positioning is deliberately complementary to specialist resources. Repbase and Dfam provide substantially deeper repeat-family references and models \citep{Bao2015Repbase,Wheeler2013Dfam}. dbRIP and euL1db address insertion polymorphisms and sample-specific L1HS records at a resolution not provided by TE-KG \citep{Wang2006dbRIP,Mir2015euL1db}. HervD Atlas offers manually curated depth for HERV-disease associations, and TE-SCALE provides single-cell cancer expression and co-expression at a scale beyond the bulk datasets described here \citep{Li2024HervDAtlas,Meng2026TESCALE}. TE-KG is therefore most useful when a question crosses evidence types, not when a user needs the deepest available record within one specialist domain.

The literature graph offers a clear provenance advantage and a clear limitation. Retaining predicates and PMIDs makes graph edges inspectable and allows users to return to the source publication. However, the current corpus was produced through a combination of automated screening, constrained language-model extraction, normalization and manual curation. Errors can arise from retrieval, abbreviation resolution, entity normalization or interpretation of the source text. PMID presence establishes traceability, not correctness or causality. A reproducible manual audit with preserved sample identifiers is required before the extraction pipeline can support a quantitative accuracy statement.

Expression and co-expression require similar caution. TE expression estimates are sensitive to repetitive sequence and quantification choices \citep{Lanciano2020TEExpression}. The three TE-KG contexts originate from different studies and cannot be interpreted as a matched normal-to-cancer experiment. The normal-tissue metadata also contain duplicated identifiers and five unmatched matrix runs. Co-expression further reduces the data to thresholded pairwise correlations and modules. FDR control limits the expected proportion of false discoveries among retained statistical tests, but it does not establish a regulatory interaction. The display catalogue adds another selection layer because it presents bounded, approved subgraphs rather than every offline edge.

The intelligent interface lowers the effort needed to retrieve evidence across several stores, but it should remain secondary to source-visible database access. Session-scoped follow-up makes short research conversations possible without building a persistent user profile. At the same time, answer quality depends on entity resolution, plugin selection, the availability of local evidence and the language model's synthesis. The interface can omit relevant literature or produce an incomplete answer even when the underlying resource contains useful records. Keeping PubMed links visible and translating uncertainty into reader-facing language are therefore more important than presenting internal reasoning steps.

Several requirements remain before submission and public release. The exact RepeatMasker and RepBase releases, acquisition dates and redistribution terms must be recorded. The expression accession-to-sample manifest must be frozen, including the SRP013565 subset and the five normal-tissue metadata mismatches. A stable public URL, versioned code and data deposits, licences and an update policy are also needed. The current \emph{Database} instructions require a free, login-free resource that remains available for at least two years, so a local deployment and Download page alone are insufficient.

The most informative next evaluation is a reproducible cross-layer use case rather than a larger catalogue of interface screenshots. A locked L1HS example can test whether a user can move from identity to classification, representative sequence or location, expression, co-expression and literature evidence without losing provenance or overstating the result. A second case involving a HERV or DNA transposon would help determine whether the workflow generalizes beyond the best represented LINE example. These cases should report missing evidence as part of the result rather than supplementing it with unverified generated text.
